% ======================================================================
% Ontology of Continua -- Core 1.3.3
% Cross-K module: crossk_k4_k5.tex
% Cross-level relations between K4 and K5
% ======================================================================

\subsubsection{Межуровневая структура K4–K5}
\label{sec:crossk-k4-k5}

Переход $K_4 \rightarrow K_5$ обозначает возникновение нейронного
континуума из биологического. Если в $K_4$ существуют градиенты,
мембраны и метаболические циклы, то в $K_5$ вводятся:
\begin{itemize}
  \item электрическая возбудимость и динамика мембранного потенциала,
  \item ионные каналы с дискретными состояниями открыто/закрыто,
  \item распространяющиеся электрические события (протоспайки),
  \item рефрактерные механизмы и пороги возбуждения,
  \item пространственное распространение вдоль мембран и ранняя сеточная структура,
  \item аттрактороподобные паттерны, формирующие протоинформационные потоки.
\end{itemize}

Это рождение организации типа нервной системы: структуры, способной
кодировать, распространять и преобразовывать электрические сигналы.

% ----------------------------------------------------------------------
\subsubsection{1. Задействованные уровни и их роли}

\paragraph{K4 (биологический континуум).}
$K_4$ предоставляет:
\begin{itemize}
  \item мембраны и компартменты $\partial\Omega$,
  \item градиенты (ионные, pH, редокс) с потенциалами $P_{\mathrm{grad}}$,
  \item ионные и осмотические потоки ($J_{\mathrm{ion}}, J_{\mathrm{osm}}$),
  \item метаболические и насосные циклы $C_{\mathrm{energy}}$,
  \item пороги $\Theta_{\mathrm{mem}}, \Theta_{\mathrm{grad}},
        \Theta_{\mathrm{osm}}, \Theta_{\mathrm{redox}}$.
\end{itemize}

Однако в $K_4$ отсутствуют:
\begin{itemize}
  \item регенеративное электрическое возбуждение,
  \item дискретная динамика открытия каналов,
  \item дальнодействующее электрическое распространение,
  \item устойчивые электрические аттракторы.
\end{itemize}

\paragraph{K5 (нейронный континуум).}
$K_5$ вводит:
\begin{itemize}
  \item ось возбудимости $A_{\mathrm{exc}}$, задающую мембранный потенциал,
  \item ионно-канальные оси $A_{\mathrm{channel}}$, описывающие состояния
        «открыто/закрыто»,
  \item электрические потенциалы $P_{\mathrm{elect}}$,
  \item потоки $J_5 = (J_{\mathrm{ion}}, J_{\mathrm{exc}}, J_{\mathrm{leak}},
        J_{\mathrm{shunt}})$,
  \item протоспайковые и рефрактерные циклы
        $C_{\mathrm{spike}},C_{\mathrm{recovery}}$,
  \item пороги $\Theta_{\mathrm{exc}}, \Theta_{\mathrm{refrac}},
        \Theta_{\mathrm{noise-electrical}}$.
\end{itemize}

$K_5$ — первый континуум, способный к пространственно-временной
паттернной передаче.

% ----------------------------------------------------------------------
\subsubsection{2. Совместные и наследуемые оси и пороги}

\paragraph{Наследуемые оси.}
Градиенты $K_4$ становятся компонентами электрической оси:
\[
A_{\mathrm{grad}}(K_4) \subset A_{\mathrm{exc}}(K_5).
\]
Мембранная ось $A_{\mathrm{mem}}$ становится подложкой для размещения
каналов и пространственного распространения.

\paragraph{Новые оси.}
$K_5$ вводит:
\begin{itemize}
  \item $A_{\mathrm{exc}}$ — ось мембранного напряжения,
  \item $A_{\mathrm{channel}}$ — дискретные состояния гейтинга,
  \item $A_{\mathrm{lat}}$ — латеральное распространение вдоль поверхности мембраны.
\end{itemize}

\paragraph{Пороги.}
Новые пороговые семейства включают:
\begin{itemize}
  \item $\Theta_{\mathrm{exc}}$ — минимальная деполяризация для формирования
        спайка,
  \item $\Theta_{\mathrm{refrac}}$ — ограничения времени восстановления,
  \item $\Theta_{\mathrm{channel}}$ — устойчивость гейтинга и точки отказа,
  \item $\Theta_{\mathrm{noise-electrical}}$ — ограничения коллапса из-за шума.
\end{itemize}

Нарушение $\Theta_4$ блокирует возбудимость:
\[
\Theta_4^{\mathrm{death}} \Rightarrow \Omega(K_5)=\varnothing.
\]

% ----------------------------------------------------------------------
\subsubsection{3. Межуровневые потоки, циклы и напряжения}

\paragraph{Потоки.}
Потоки $K_5$ наследуются от $J_4$, но получают новые электрические
компоненты:
\[
J_5 = (J_{\mathrm{ion}}, J_{\mathrm{exc}}, J_{\mathrm{leak}}, J_{\mathrm{shunt}}).
\]

Ключевые вкладки:
\begin{itemize}
  \item \textbf{$J_{\mathrm{exc}}$}: регенеративный натриевый-like всплеск,
  \item \textbf{$J_{\mathrm{leak}}$}: затухающие фоновый ток,
  \item \textbf{$J_{\mathrm{shunt}}$}: мембранные стабилизирующие обходные токи,
  \item \textbf{$J_{\mathrm{ion}}$}: классический градиентный поток по типу
        уравнения Нернста.
\end{itemize}

\paragraph{Циклы.}
В $K_5$ появляются два фундаментальных цикла:
\[
C_{\mathrm{spike}}:\ \Delta V \uparrow \rightarrow \Delta V_{\mathrm{peak}}
\rightarrow \text{реполяризация} \rightarrow \text{восстановление}.
\]

\[
C_{\mathrm{recovery}}:\ \text{рефрактерность} \rightarrow \text{сброс}
\rightarrow \text{восстановленная возбудимость}.
\]

Эти циклы зависят от устойчивых градиентов и кинетики каналов.

\paragraph{Напряжение.}
Межуровневое напряжение:
\[
T_{4,5} =
g(\Theta_{\mathrm{exc}} - P_{\mathrm{elect}},\
  J_{\mathrm{leak}} - J_{\mathrm{exc}},\
  \Theta_{\mathrm{channel}} - A_{\mathrm{channel}}).
\]

Когда $T_{4,5}$ превышает размерностный порог, распространение прекращается,
и система возвращается к неэксцитируемому режиму K₄.

% ----------------------------------------------------------------------
\subsubsection{4. Условия рождения и исчезновения между уровнями}

\paragraph{Рождение \texorpdfstring{$K_5$}{K_5}.}
Нейронный континуум возникает, когда:
\[
T_4 > \Theta_{4,\mathrm{dim}}
\quad\text{и}\quad
A_{\mathrm{exc}}, A_{\mathrm{channel}} \in M_5 \setminus A(K_4),
\]
что позволяет устойчиво развивать спайковую динамику.

Расширение пространства состояний:
\[
\Omega(K_5) = \Omega(K_4) \cup \Delta\Omega_{\mathrm{neural}}.
\]

\paragraph{Распространение гибели.}
Если пороги $\Theta_4$ не выполняются (разрыв мембраны, схлопывание градиентов),
то $K_5$ теряет основу для возбудимости.

Если пороги $\Theta_5$ не выполняются (избыточный шум, нестабильность каналов,
недостаточный градиент), разрушается только $K_5$; $K_4$ выживает.

% ----------------------------------------------------------------------
\subsubsection{5. Концептуальные примеры}

\paragraph{Формирование протоспайка.}
При выполнении:
\[
\Delta V > \Theta_{\mathrm{exc}}
\quad\text{и}\quad
A_{\mathrm{channel}} = \text{open},
\]
воспроизводится регенеративный подъем и распространение по локальной мембране.

\paragraph{2D-перенос.}
Латеральная ось $A_{\mathrm{lat}}$ позволяет распространять электрические
паттерны по двумерным поверхностям мембраны — это невозможно в K₄.

\paragraph{Возбудимость как новая размерность.}
Появление динамики $\Delta V(t)$ составляет новую ось и новый класс потоков,
удовлетворяя правилу размерностного роста:
\[
\dim K_5 = \dim K_4 + 1.
\]

\paragraph{Ранняя информационная передача.}
Спайкоподобная активность поддерживает передачу паттернов и протокодирование,
создавая первый субстрат для структурированной информации в иерархии континуумов.

Следовательно, переход K4→K5 можно сжато описать как рождение
возбудимости, электрической сигнализации и распространения паттернов из
биологических градиентов и мембран.
